\documentclass[11pt]{article}

\usepackage[T1]{fontenc}
\usepackage[utf8]{inputenc}
\usepackage{lmodern}
\usepackage[a4paper,margin=27mm]{geometry}
\usepackage{amsmath,amssymb,amsthm,mathtools}
\usepackage{microtype}
\usepackage{xcolor}
\usepackage{enumitem}
\usepackage{hyperref}

\hypersetup{
  colorlinks=true,
  linkcolor=blue!55!black,
  citecolor=blue!55!black,
  urlcolor=blue!55!black
}

\newtheorem{theorem}{Theorem}[section]
\newtheorem{proposition}[theorem]{Proposition}
\newtheorem{lemma}[theorem]{Lemma}
\newtheorem{corollary}[theorem]{Corollary}
\newtheorem{definition}[theorem]{Definition}
\theoremstyle{remark}
\newtheorem{remark}[theorem]{Remark}

\newcommand{\CAR}{\operatorname{CAR}}
\newcommand{\HS}{\mathrm{HS}}
\newcommand{\Tr}{\operatorname{Tr}}
\newcommand{\id}{\mathbf 1}
\newcommand{\Z}{\mathbb Z}
\newcommand{\C}{\mathbb C}
\newcommand{\A}{\mathcal A}
\newcommand{\F}{\mathcal F}
\newcommand{\K}{\mathcal K}
\newcommand{\norm}[1]{\left\lVert #1\right\rVert}
\newcommand{\ip}[2]{\left\langle #1,#2\right\rangle}

\title{\texorpdfstring{$\Psi_\lambda$}{Psi-lambda} at the Quasi-Free Boundary:\\
Finite-Window Implementers, the \texorpdfstring{$\mathbb Z_4$}{Z4} Tower,
and the Two-Edge Shale Obstruction}
\author{Proof memorandum}
\date{28 August 2026}

\begin{document}
\maketitle

\begin{abstract}
We prove the strongest statement presently justified by the frozen
finite-collar data on the proposed field $\Psi_\lambda$.  On every fixed
finite window, trace-norm convergence of the restricted covariances gives,
after the canonical Araki--Wyss purification, canonical Bogoliubov
implementers with an explicit Shale bound.  Along every coherent
subsequence these implementers converge, in fact in operator norm; the
estimate on arbitrary finite-particle vectors is written out.  For a
compatible inner $\mathbb Z_4$ tower we prove directly that crossed product
commutes with the inductive limit and that the canonical expectations are
unique with their necessary normalisations.  Finally, for the exact
two-edge chiral free-fermion model we compute the Hilbert--Schmidt
divergence caused by a quarter twist:
\[
 \norm{P_0-P_{1/4}}_{\HS}^{\,2}
 =\frac{2}{\pi^2}\log N+O(1),
 \qquad \frac{2}{\pi^2}=0.202642367\ldots .
\]
This agrees to $0.18\%$ with the measured QWZ-collar slope $0.203$.
The document also proves that the measurements stated in the repository
do \emph{not} establish the stronger net-level conclusion: compressed
covariance data do not determine a Bogoliubov transformation, four energy
moments do not prove an all-orders energy bound, and a finite logarithmic
fit does not prove an asymptotic law.  The missing net-identification
statement is isolated verbatim at the end.
\end{abstract}

\tableofcontents

\section{Scope, frozen input, and the necessary correction}

\subsection{The frozen input}

The one-particle lattice Hamiltonian used in the quantitative probes is
the Qi--Wu--Zhang collar
\[
 h(k_x,k_y)=\sin k_x\,\sigma_x+\sin k_y\,\sigma_y+
 (M-\cos k_x-\cos k_y)\sigma_z
\]
on a cylinder of circumference $N_x$, open width $N_y=8$, and boundary
twist $i^k$, $k=0,1,2,3$.  The topological and trivial controls are
$M=1$ and $M=3$, respectively.  If $H_N^{(k)}$ denotes the finite
one-particle matrix, the particle-hole symmetric covariance is
\[
 C_N^{(k)}
 =\chi_{(-\infty,0)}(H_N^{(k)})
   +\frac12\chi_{\{0\}}(H_N^{(k)}).
\]
The equal-particle Slater projection, used when the energy-sign fillings
have different particle number, is denoted by $P_N^{(k)}$.

The following numbers are inputs to this memorandum, not conclusions of
the proof:
\begin{enumerate}[label=(D\arabic*)]
\item On fixed edge windows of widths $6$ and $8$, the compressed
  sector differences are Cauchy, with a representative fitted rate
  approximately $N_x^{-1.85}$.  The one-sided quantity
  $\norm{(C_N^{(0)}-C_N^{(1)})E_w}_{\HS}^2$ plateaus below $0.376$.
\item The restricted covariances converge in trace norm on every tested
  fixed window.
\item In the topological run,
  \[
  \norm{C_N^{(0)}-C_N^{(1)}}_{\HS}^2
  =3.605+0.203\log N_x
  \]
  is the reported fit, while the trivial control saturates near $0.546$.
\item On one polar-phase-fixed Pfaffian band, selected matrix elements
  are Cauchy.  A second band has a different limit and is not mixed with
  the first.
\item For $n=1,\ldots,4$, the energy ratios tend toward $1$, and the
  measured excitation norms obey
  $f_n<C^{n+1}(n!)^2$ with probe-provenance constant $C=2.86$.
\item The exact finite data give
  \[
  \lambda=(\omega_s,\omega_f),\qquad
  \norm{\lambda}^2=2,\qquad h_\lambda=1,\qquad
  [\lambda]^4=[0],\qquad [\lambda]^2=[v].
  \]
  The finite $\mathbb Z_4$ clock action is compatible with the tested
  tower embeddings.
\end{enumerate}

\subsection{Why the originally proposed implication is false}

Three distinctions are mathematically indispensable.

\paragraph{A compressed covariance is not a polarization.}
If $E$ is a proper window projection and $P$ is a global Fermi projection,
then $EPE$ is generally a positive contraction, not a projection:
\[
 (EPE)^2-EPE=-EP(\id-E)PE.
\]
Consequently, Shale--Stinespring cannot be applied to $EPE$ as though it
were a Fock polarization.  We repair this by the canonical
Araki--Wyss purification in Section~\ref{sec:local}.

\paragraph{Two covariances do not specify a sector automorphism.}
Even in finite dimension, two mixed covariances $C$ and $D$ need not be
unitarily conjugate: unitary conjugacy preserves the spectrum.  Thus
$\norm{C-D}_{1}<\infty$ or $C_N-D_N\to0$ does not by itself produce a
Bogoliubov map on the original window CAR algebra.  It does produce a
canonical equivalence of the purified standard forms, which is the exact
local result proved below.

\paragraph{Finite observations are not asymptotic theorems.}
A fit on $16\le N_x\le96$ cannot prove divergence as $N_x\to\infty$, and
inequalities checked only for $n\le4$ cannot prove an energy estimate for
every $n$.  We therefore prove the logarithm exactly for the universal
two-chiral-edge model and state the precise uniform-remainder hypothesis
under which its coefficient transfers to a lattice collar.

\section{Classical theorems used, with full statements}

Only two external analytic results are used in the proof.  They are
quoted here in the form needed below.

\begin{theorem}[Shale--Stinespring criterion for CAR]\label{thm:ss}
Let $\K_{\mathbb R}$ be a real Hilbert space, let $J$ be an orthogonal
complex structure ($J^*=-J$ and $J^2=-\id$), and let $\pi_J$ be the
irreducible fermionic Fock representation of the Clifford algebra
determined by $J$.  If $R$ is an orthogonal operator on
$\K_{\mathbb R}$, then there exists a unitary $U_R$ on the Fock space
such that
\[
 U_R\pi_J(B(f))U_R^*=\pi_J(B(Rf))
 \quad\text{for every }f\in\K_{\mathbb R}
\]
if and only if $[R,J]$ is Hilbert--Schmidt.  When it exists, $U_R$ is
unique up to multiplication by a complex scalar of modulus one.
\end{theorem}

This is the CAR theorem of D.~Shale and W.~F.~Stinespring,
\emph{States of the Clifford algebra}, Annals of Mathematics
\textbf{80} (1964), 365--381, in particular the implementability
criterion in the main theorem on pp.~375--376.  In the equivalent
polarization notation, with $P=(\id-iJ)/2$ and $Q=RPR^*$, the criterion
is
\[
 [P,R]\in\mathcal L^2
 \quad\Longleftrightarrow\quad
 P-Q\in\mathcal L^2,
\]
because
\[
 [P,R]R^*=P-RPR^*=P-Q.
\]
Thus the two Hilbert--Schmidt norms are equal.  We will verify the
hypothesis by an explicit estimate.

\begin{theorem}[Powers--St{\o}rmer square-root inequality]
\label{thm:ps}
Let $A$ and $B$ be positive trace-class operators on a Hilbert space.
Then
\[
 \norm{A^{1/2}-B^{1/2}}_2^2\le \norm{A-B}_1.
\]
The same assertion applies to $\id-A$ and $\id-B$ whenever
$0\le A,B\le\id$ and the differences in question are trace class.
\end{theorem}

This is the square-root estimate proved in
R.~T.~Powers and E.~St{\o}rmer,
\emph{Free states of the canonical anticommutation relations},
Communications in Mathematical Physics \textbf{16} (1970), 1--33.
The finite-dimensional application below satisfies all hypotheses:
the covariance matrices and their complements are positive, and every
operator on the fixed window is trace class.

\begin{remark}[The pure-state quasi-equivalence form]
H.~Araki, \emph{On quasifree states of CAR and Bogoliubov
automorphisms}, Publications of the Research Institute for Mathematical
Sciences \textbf{6} (1970/71), 385--442, proves the general
quasi-equivalence criterion.  For pure gauge-invariant quasi-free states
with polarizations $P$ and $Q$, it reduces to:
\[
 \pi_P\text{ and }\pi_Q\text{ are unitarily equivalent}
 \quad\Longleftrightarrow\quad P-Q\in\mathcal L^2.
\]
We use this form only to interpret the global obstruction.  The local
construction itself is given explicitly and does not defer to Araki's
theorem.
\end{remark}

\section{Finite-window implementers and convergence}
\label{sec:local}

\subsection{Canonical purification of a window covariance}

Let $\mathfrak h_w$ be the finite-dimensional one-particle space of a
fixed window and let $C$ be a covariance, so $0\le C\le\id$.  Put
\[
 V_C:\mathfrak h_w\longrightarrow\mathfrak h_w\oplus\mathfrak h_w,
 \qquad
 V_C f=
 \begin{pmatrix}C^{1/2}f\\(\id-C)^{1/2}f\end{pmatrix},
\]
and define
\[
 \widehat P_C=V_CV_C^*
 =
 \begin{pmatrix}
 C& C^{1/2}(\id-C)^{1/2}\\
 C^{1/2}(\id-C)^{1/2}&\id-C
 \end{pmatrix}.
\]

\begin{lemma}[Purification and its quantitative continuity]
\label{lem:purification}
For every covariance $C$, $\widehat P_C$ is an orthogonal projection of
rank $\dim\mathfrak h_w$.  For any two covariances $C,D$,
\[
 \norm{\widehat P_C-\widehat P_D}_2
 \le 2\sqrt{2\norm{C-D}_1}.
\tag{3.1}\label{eq:purification-bound}
\]
\end{lemma}

\begin{proof}
Functional calculus makes $C^{1/2}$ commute with
$(\id-C)^{1/2}$, and
\[
 V_C^*V_C=C+(\id-C)=\id.
\]
Hence $V_C$ is an isometry.  Therefore
\[
 \widehat P_C^*=\widehat P_C,\qquad
 \widehat P_C^2
 =V_C(V_C^*V_C)V_C^*
 =V_CV_C^*=\widehat P_C.
\]
Its range is the range of the isometry $V_C$, so its rank is
$\dim\mathfrak h_w$.

Next,
\[
\begin{aligned}
\widehat P_C-\widehat P_D
 &=(V_C-V_D)V_C^*+V_D(V_C^*-V_D^*).
\end{aligned}
\]
The ideal property
$\norm{XY}_2\le\norm{X}_2\norm{Y}$ and
$\norm{V_C}=\norm{V_D}=1$ give
\[
 \norm{\widehat P_C-\widehat P_D}_2
 \le 2\norm{V_C-V_D}_2.
\]
Moreover,
\[
\begin{aligned}
\norm{V_C-V_D}_2^2
 &=
 \norm{C^{1/2}-D^{1/2}}_2^2\\
 &\quad+
 \norm{(\id-C)^{1/2}-(\id-D)^{1/2}}_2^2\\
 &\le \norm{C-D}_1+\norm{(\id-C)-(\id-D)}_1\\
 &=2\norm{C-D}_1,
\end{aligned}
\]
where Theorem~\ref{thm:ps} was used twice.  Combining the last two
inequalities proves \eqref{eq:purification-bound}.
\end{proof}

\subsection{The direct rotation}

\begin{lemma}[Canonical direct rotation]\label{lem:rotation}
Let $P,Q$ be finite-rank orthogonal projections of equal rank and suppose
$\norm{P-Q}<1$.  Set
\[
 S=QP+(\id-Q)(\id-P),\qquad
 R=S(S^*S)^{-1/2}.
\tag{3.2}\label{eq:direct-rotation}
\]
Then $R$ is unitary, $RPR^*=Q$, and $R$ depends continuously in operator
norm on $(P,Q)$ throughout the open set $\norm{P-Q}<1$.  In addition,
\[
 \norm{[P,R]}_2=\norm{P-Q}_2.
\tag{3.3}\label{eq:shale-equality}
\]
\end{lemma}

\begin{proof}
First,
\[
 SP=QP=QS.
\tag{3.4}\label{eq:intertwine-S}
\]
Taking adjoints gives $PS^*=S^*Q$.  Hence
\[
 P(S^*S)=S^*QS=(S^*S)P,
\]
so $P$ commutes with every continuous function of $S^*S$, including
$(S^*S)^{-1/2}$ once invertibility is established.

For $x=Px+(\id-P)x$, the two terms in $Sx$ lie in the orthogonal
subspaces $Q\K$ and $(\id-Q)\K$, and therefore
\[
 \norm{Sx}^2
 =\norm{QPx}^2+\norm{(\id-Q)(\id-P)x}^2.
\]
For $y\in P\K$,
\[
 \norm{Qy}\ge \norm{y}-\norm{(\id-Q)y}
 \ge (1-\norm{P-Q})\norm{y},
\]
because $(\id-Q)y=(P-Q)y$.  The same argument on $(\id-P)\K$ gives
\[
 \norm{(\id-Q)z}\ge(1-\norm{P-Q})\norm{z}.
\]
Thus
\[
 \norm{Sx}\ge(1-\norm{P-Q})\norm{x}.
\]
So $S$ is injective and bounded below.  In finite dimension it is
invertible; hence $S^*S$ is strictly positive and
\eqref{eq:direct-rotation} is defined.

The polar factor $R$ is unitary:
\[
 R^*R=(S^*S)^{-1/2}S^*S(S^*S)^{-1/2}=\id,
\]
and, since $S$ is invertible in equal finite dimensions, also $RR^*=\id$.
Using \eqref{eq:intertwine-S} and commutation of $P$ with
$(S^*S)^{-1/2}$,
\[
 RP=S(S^*S)^{-1/2}P
 =SP(S^*S)^{-1/2}
 =QS(S^*S)^{-1/2}=QR.
\]
Multiplication by $R^*$ yields $RPR^*=Q$.

All operations in \eqref{eq:direct-rotation} are norm-continuous where
$S^*S$ is bounded away from zero.  The displayed lower bound supplies
that separation locally, proving continuity.  Finally,
\[
 [P,R]R^*=P-RPR^*=P-Q.
\]
Right multiplication by a unitary preserves the Hilbert--Schmidt norm,
which proves \eqref{eq:shale-equality}.
\end{proof}

\begin{theorem}[Complete finite-window theorem]\label{thm:local}
Fix a finite window $\mathfrak h_w$.  For $j\in\mathbb N$ and
$r\in\{0,1\}$ let $C_j^{(r)}$ be covariances on $\mathfrak h_w$.
Assume
\[
 C_j^{(r)}\longrightarrow C^{(r)}
 \quad\text{in trace norm for }r=0,1
\tag{3.5}\label{eq:trace-conv}
\]
and
\[
 \norm{\widehat P_{C^{(0)}}-\widehat P_{C^{(1)}}}<1.
\tag{3.6}\label{eq:angle}
\]
Then, after discarding finitely many $j$, the following statements hold.

\begin{enumerate}[label=\textup{(\alph*)}]
\item The direct rotation $R_j$ from
  $\widehat P_{C_j^{(0)}}$ to $\widehat P_{C_j^{(1)}}$ is defined.
  Its Bogoliubov automorphism in the finite-window Araki--Wyss standard
  representation has a unitary implementer $U_j$, unique up to phase.
  The exact Shale number and a covariance-only bound are
  \[
  \norm{[\widehat P_{C_j^{(0)}},R_j]}_2
  =\norm{\widehat P_{C_j^{(0)}}-\widehat P_{C_j^{(1)}}}_2
  \le2\sqrt{2\norm{C_j^{(0)}-C_j^{(1)}}_1}.
  \tag{3.7}\label{eq:local-shale}
  \]
\item $R_j\to R$ in operator norm, where $R$ is the direct rotation
  of the limiting purified projections.  The phases may be chosen
  coherently so that $U_j\to U$ in operator norm, and hence strongly,
  on the full finite-window Fock space.
\item If in addition $C^{(0)}=C^{(1)}$, then $R=\id$, the phase may be
  fixed by $\ip{\Omega}{U_j\Omega}>0$ for all sufficiently large $j$,
  and $U_j\to\id$.
\end{enumerate}
\end{theorem}

\begin{proof}
By Lemma~\ref{lem:purification} and \eqref{eq:trace-conv},
\[
 \widehat P_{C_j^{(r)}}\longrightarrow\widehat P_{C^{(r)}}
 \quad\text{in Hilbert--Schmidt norm, hence in operator norm}.
\]
Condition \eqref{eq:angle} is strict, so for all sufficiently large $j$,
\[
 \norm{\widehat P_{C_j^{(0)}}-\widehat P_{C_j^{(1)}}}<1.
\]
Lemma~\ref{lem:rotation} defines $R_j$ and proves
\eqref{eq:shale-equality}.  Lemma~\ref{lem:purification}, applied to
$C_j^{(0)}$ and $C_j^{(1)}$, gives the last inequality in
\eqref{eq:local-shale}.  This verifies the Hilbert--Schmidt hypothesis of
Theorem~\ref{thm:ss}; in finite dimension it is of course automatically
finite, but \eqref{eq:local-shale} is the promised quantitative bound.
Uniqueness up to phase follows either from Theorem~\ref{thm:ss} or
directly: if $U_j$ and $U'_j$ implement the same automorphism, then
$U_j'^*U_j$ commutes with every represented CAR generator, and the
finite Fock representation is irreducible, so $U_j'^*U_j$ is scalar.

Continuity in Lemma~\ref{lem:rotation} gives $R_j\to R$.  We now prove,
rather than merely cite, the implementer continuity.  Regard the doubled
complex one-particle space as a real Euclidean space of dimension $2d$,
and choose Majorana generators
$\gamma_1,\ldots,\gamma_{2d}$ satisfying
\[
 \gamma_a^*=\gamma_a,\qquad
 \gamma_a\gamma_b+\gamma_b\gamma_a=2\delta_{ab}\id .
\]
For all sufficiently large $j$, $R^*R_j$ lies in the logarithm
neighbourhood of the identity.  Let
\[
 A_j=\log(R^*R_j);
\]
then $A_j$ is real skew-adjoint and $\norm{A_j}\to0$.  Define
\[
 K_j=\frac14\sum_{a,b=1}^{2d}(A_j)_{ab}\gamma_a\gamma_b.
\]
Since the sum is finite,
\[
 \norm{K_j}
 \le\frac14\sum_{a,b}|(A_j)_{ab}|
 \le \frac{d}{2}\norm{A_j}_2
 \longrightarrow0.
\tag{3.8}\label{eq:spin-bound}
\]
The Clifford commutator is computed without omission:
\[
\begin{aligned}
[\gamma_a\gamma_b,\gamma_c]
&=\gamma_a\gamma_b\gamma_c-\gamma_c\gamma_a\gamma_b\\
&=\gamma_a(-\gamma_c\gamma_b+2\delta_{bc})
 -(-\gamma_a\gamma_c+2\delta_{ac})\gamma_b\\
&=2\delta_{bc}\gamma_a-2\delta_{ac}\gamma_b.
\end{aligned}
\]
Consequently,
\[
 [K_j,\gamma_c]
 =\sum_a(A_j)_{ac}\gamma_a.
\]
Solving the finite-dimensional differential equation
$\frac d{dt}(e^{tK_j}\gamma_ce^{-tK_j})
 =[K_j,e^{tK_j}\gamma_ce^{-tK_j}]$
shows that $e^{K_j}$ implements $e^{A_j}=R^*R_j$.
If $U$ is one implementer of $R$, choose
\[
 U_j=Ue^{K_j}.
\]
Then $U_j$ implements $R_j$, and
\[
\begin{aligned}
\norm{U_j-U}
 &=\norm{e^{K_j}-\id}\\
 &=\norm{\int_0^1 e^{tK_j}K_j\,dt}\\
 &\le\int_0^1e^{t\norm{K_j}}\norm{K_j}\,dt\\
 &\le e^{\norm{K_j}}\norm{K_j}\longrightarrow0.
\end{aligned}
\tag{3.9}\label{eq:implementer-norm}
\]
This proves operator-norm convergence.

For completeness, here is the requested estimate on an arbitrary
$n$-particle core vector.  Let
\[
 \xi=a^*(f_1)\cdots a^*(f_n)\Omega .
\]
Write $\beta_j(a^*(f))=U_ja^*(f)U_j^*$ and
$\beta(a^*(f))=Ua^*(f)U^*$.  A Bogoliubov transform is linear in the
creation and annihilation parts, so
\[
 \norm{\beta_j(a^*(f))-\beta(a^*(f))}
 \le\norm{R_j-R}\,\norm{f}.
\tag{3.10}\label{eq:field-bound}
\]
Insert and subtract one factor at a time:
\[
\begin{aligned}
U_j\xi-U\xi
={}&
\beta_j(a^*(f_1))\cdots\beta_j(a^*(f_n))
  (U_j\Omega-U\Omega)\\
&+\sum_{\ell=1}^n
\beta_j(a^*(f_1))\cdots\beta_j(a^*(f_{\ell-1}))\\
&\qquad\cdot
\bigl(\beta_j(a^*(f_\ell))-\beta(a^*(f_\ell))\bigr)
\beta(a^*(f_{\ell+1}))\cdots\beta(a^*(f_n))U\Omega .
\end{aligned}
\]
Because $\norm{a^*(f)}=\norm f$ and automorphisms preserve norms,
\[
\norm{U_j\xi-U\xi}
\le
\left(
\norm{U_j\Omega-U\Omega}
n\norm{R_j-R}
\right)\prod_{\ell=1}^n\norm{f_\ell}.
\tag{3.11}\label{eq:nparticle}
\]
Both terms in parentheses tend to zero by
\eqref{eq:implementer-norm} and $R_j\to R$.  Finite linear combinations
of such vectors are handled by the triangle inequality.  This gives the
strong-convergence argument explicitly on the finite-particle core.

If $C^{(0)}=C^{(1)}$, the two limiting purified projections coincide.
Formula \eqref{eq:direct-rotation} then has $S=\id$ and $R=\id$.
Equation \eqref{eq:implementer-norm} permits $U=\id$.  Since
$U_j\Omega\to\Omega$, the overlap is nonzero for large $j$; multiplying
$U_j$ by its unique phase makes the overlap positive, without affecting
convergence.
\end{proof}

\begin{corollary}[Bound in terms of a measured window HS norm]
\label{cor:measured}
If $m=\dim\mathfrak h_w$, then
\[
 \norm{C-D}_1\le\sqrt m\,\norm{C-D}_2,
\]
and therefore
\[
 \norm{[\widehat P_C,R]}_2
 \le2\sqrt2\,m^{1/4}\norm{C-D}_2^{1/2}.
\tag{3.12}
\]
\end{corollary}

\begin{proof}
If $s_1,\ldots,s_m$ are the singular values of $C-D$, Cauchy--Schwarz
gives
\[
 \norm{C-D}_1=\sum_{\ell=1}^m s_\ell
\le\sqrt m\left(\sum_{\ell=1}^m s_\ell^2\right)^{1/2}
=\sqrt m\,\norm{C-D}_2.
\]
Substitution into \eqref{eq:local-shale} proves the result.
\end{proof}

\begin{remark}[Exactly what the window theorem proves]
Theorem~\ref{thm:local} proves implementability in the
\emph{purified finite-window standard form}.  It does not assert that the
two mixed restricted states are related by an automorphism of
$\CAR(\mathfrak h_w)$ itself.  Such an assertion would require, at a
minimum, equality of the spectra of the two restricted covariances.
Nor does the theorem construct a single nontrivial global charged field:
when the sectors collapse locally, the canonical local direct rotations
converge to the identity.  The charge is detected by the failure to
assemble these local equivalences into a global implementer.
\end{remark}

\section{The exact \texorpdfstring{$\mathbb Z_4$}{Z4} tower}
\label{sec:tower}

Let
\[
 \A_1\xrightarrow{\iota_1}\A_2\xrightarrow{\iota_2}\cdots
\]
be a unital inductive system of finite matrix algebras.  Suppose
$u_N\in\A_N$ is unitary,
\[
 u_N^4=\id,\qquad \iota_N(u_N)=u_{N+1},
\tag{4.1}\label{eq:compatible-clock}
\]
and let $\alpha_N=\operatorname{Ad}(u_N)$.
Write $v_N$ for the canonical crossed-product unitary:
\[
 v_Nav_N^*=\alpha_N(a),\qquad v_N^4=\id.
\]

\begin{theorem}[Naturality, fusion, and expectations]
\label{thm:crossed}
Under \eqref{eq:compatible-clock}:
\begin{enumerate}[label=\textup{(\alph*)}]
\item There is a natural $*$-isomorphism
\[
 \Phi_N:\A_N\rtimes_{\alpha_N}\mathbb Z_4
 \longrightarrow
 \A_N\otimes C^*(\mathbb Z_4),
\qquad
 \Phi_N(av_N^g)=au_N^g\otimes w^g,
\tag{4.2}\label{eq:phi}
\]
where $w$ is the canonical generator of $C^*(\mathbb Z_4)$.
\item If $\A=\varinjlim\A_N$ and $\alpha=\varinjlim\alpha_N$, then
\[
 \varinjlim(\A_N\rtimes_{\alpha_N}\mathbb Z_4)
 \cong\A\rtimes_\alpha\mathbb Z_4
 \cong\A\otimes C^*(\mathbb Z_4).
\tag{4.3}\label{eq:limit-crossed}
\]
\item The homogeneous subspaces
$\mathcal B_g=\overline{\A v^g}$ obey
\[
 \mathcal B_g\mathcal B_h\subset\mathcal B_{g+h},
\qquad
\mathcal B_g^*=\mathcal B_{-g},
\qquad
v^4=\id\in\A.
\tag{4.4}\label{eq:fusion}
\]
Thus a degree-one symbol $\Psi_\lambda:=v$ has
$\Psi_\lambda^4$ in the neutral algebra, and
$\Psi_\lambda^2$ is the unique order-two grade.  With the frozen label
$[\lambda]^2=[v]$, this is exactly the abstract $\mathbb Z_4$ fusion
law.
\item The coefficient map
\[
 F\left(\sum_{g=0}^3a_gv^g\right)=a_0
\tag{4.5}\label{eq:coefficient-expectation}
\]
extends to a faithful conditional expectation
$F:\A\rtimes_\alpha\mathbb Z_4\to\A$.  It is the unique conditional
expectation invariant under the dual $\widehat{\mathbb Z_4}$ action.
\item If $\tau_N$ is a compatible faithful trace on the tower, then
\[
 E_N(a)=\frac14\sum_{g=0}^3\alpha_N^g(a)
\tag{4.6}\label{eq:average-expectation}
\]
is the unique $\tau_N$-preserving conditional expectation
$\A_N\to\A_N^{\mathbb Z_4}$.  The $E_N$ are natural and induce the
corresponding trace-preserving expectation on the inductive limit.
\end{enumerate}
\end{theorem}

\begin{proof}
Map $a\in\A_N$ to $a\otimes\id$ and $v_N$ to $u_N\otimes w$.
The defining relations are preserved:
\[
\begin{aligned}
(u_N\otimes w)(a\otimes\id)(u_N^*\otimes w^*)
 &=u_Nau_N^*\otimes\id=\alpha_N(a)\otimes\id,\\
(u_N\otimes w)^4&=u_N^4\otimes w^4=\id.
\end{aligned}
\]
The universal property of the crossed product therefore gives a
$*$-homomorphism $\Phi_N$.  It is onto because its range contains
$a\otimes\id$ and also
\[
 \Phi_N(u_N^*v_N)=\id\otimes w.
\]
An inverse is given on generators by
\[
 a\otimes\id\longmapsto a,\qquad
 \id\otimes w\longmapsto u_N^*v_N.
\]
Indeed, $u_N^*v_N$ commutes with $\A_N$:
\[
 (u_N^*v_N)a
 =u_N^*\alpha_N(a)v_N
 =u_N^*u_Nau_N^*v_N
 =a(u_N^*v_N),
\]
and $(u_N^*v_N)^4=\id$ because $v_Nu_Nv_N^*=u_N$.
Thus \eqref{eq:phi} is an isomorphism.

For naturality, let
$\widetilde\iota_N:\A_N\rtimes\mathbb Z_4
\to\A_{N+1}\rtimes\mathbb Z_4$
send $a$ to $\iota_N(a)$ and $v_N$ to $v_{N+1}$.  On a spanning
monomial,
\[
\begin{aligned}
\Phi_{N+1}\widetilde\iota_N(av_N^g)
 &=\iota_N(a)u_{N+1}^g\otimes w^g\\
 &=\iota_N(au_N^g)\otimes w^g\\
 &=(\iota_N\otimes\id)\Phi_N(av_N^g).
\end{aligned}
\tag{4.7}\label{eq:naturality}
\]
So the diagrams commute exactly.  Taking inductive limits in
\eqref{eq:naturality} gives the first isomorphism in
\eqref{eq:limit-crossed}; applying the same generator argument to the
limit unitary $u=\varinjlim u_N$ gives the second.

For $a,b\in\A$,
\[
 (av^g)(bv^h)=a\alpha^g(b)v^{g+h}\in\A v^{g+h},
\]
and
\[
 (av^g)^*=v^{-g}a^*=\alpha^{-g}(a^*)v^{-g}.
\]
These identities prove \eqref{eq:fusion}, including $v^4=\id$.
The order-two identification is precisely the frozen lattice equality
$[\lambda]^2=[v]$; no net identification is inferred from the notation.

Under $\Phi$, the map $F$ is
\[
 \id_\A\otimes\tau_{\mathbb Z_4},
\qquad
\tau_{\mathbb Z_4}(w^g)=\delta_{g0}.
\]
The group trace is positive and faithful on the finite-dimensional
algebra $C^*(\mathbb Z_4)$, hence $F$ is positive, unital, faithful,
idempotent, fixes $\A$, and satisfies
$F(aXb)=aF(X)b$ for $a,b\in\A$.  Thus it is a faithful conditional
expectation.

Let $\widehat\alpha_\chi$ be the dual action,
$\widehat\alpha_\chi(a)=a$ and
$\widehat\alpha_\chi(v)=\chi(1)v$.  If $G$ is a dual-invariant
conditional expectation onto $\A$, then for $g\ne0$ choose a character
$\chi$ with $\chi(g)\ne1$.  For every $a\in\A$,
\[
 G(av^g)
 =G(\widehat\alpha_\chi(av^g))
 =\chi(g)G(av^g),
\]
so $G(av^g)=0$.  Also $G(a)=a$.  Hence $G=F$ on the algebraic crossed
product, a dense subalgebra, and continuity gives $G=F$ everywhere.

Finally, \eqref{eq:average-expectation} is plainly unital and fixes
$\A_N^{\mathbb Z_4}$.  It is positive because it is an average of
$*$-automorphisms, idempotent because
\[
 E_N(E_N(a))
 =\frac1{16}\sum_{g,h=0}^3\alpha_N^{g+h}(a)
 =\frac14\sum_{\ell=0}^3\alpha_N^\ell(a)=E_N(a),
\]
and bimodular because, for fixed $x,y$,
\[
 E_N(xay)=\frac14\sum_gx\alpha_N^g(a)y=xE_N(a)y.
\]
It preserves $\tau_N$ because traces are invariant under inner
automorphisms.

To prove uniqueness, equip $\A_N$ with
$\ip{x}{y}_{2,\tau_N}=\tau_N(x^*y)$.  For $b$ fixed by the action,
\[
\begin{aligned}
\ip{b}{a-E_N(a)}_{2,\tau_N}
&=\tau_N(b^*a)
 -\frac14\sum_g\tau_N(b^*\alpha_N^g(a))\\
&=\tau_N(b^*a)
 -\frac14\sum_g\tau_N(\alpha_N^{-g}(b^*)a)\\
&=\tau_N(b^*a)-\tau_N(b^*a)=0.
\end{aligned}
\]
Thus $E_N(a)$ is the orthogonal projection of $a$ onto the fixed
subspace.  If $G_N$ is another $\tau_N$-preserving conditional
expectation, then for fixed $b$ the bimodule property gives
\[
\ip{b}{a-G_N(a)}_{2,\tau_N}
=\tau_N(b^*a)-\tau_N(G_N(b^*a))
=\tau_N(b^*a)-\tau_N(b^*G_N(a))=0.
\]
So $G_N(a)$ is the same orthogonal projection, and $G_N=E_N$.
Compatibility follows by substituting
\eqref{eq:compatible-clock} into \eqref{eq:average-expectation}.
\end{proof}

\begin{remark}[Why the qualifier on uniqueness is necessary]
For an inner action,
$\A\rtimes\mathbb Z_4\cong\A\otimes C^*(\mathbb Z_4)$.
An arbitrary state on the second tensor factor defines a conditional
expectation onto $\A$, so a bare claim of uniqueness is false.
The coefficient expectation is unique after imposing dual invariance
(equivalently, the canonical group trace).  Separately, the averaging
expectation $\A\to\A^{\mathbb Z_4}$ is unique after imposing preservation
of the specified trace.  These are the two precise uniqueness statements
supported by the finite data.
\end{remark}

\section{The quarter-twist logarithm and global obstruction}
\label{sec:global}

\subsection{One chiral edge}

Let $\mathfrak h=L^2([0,1])$.  For $n,m\in\Z$ and
$0<\alpha<1$, define
\[
 e_n(x)=e^{2\pi inx},\qquad
 f_m^{(\alpha)}(x)=e^{2\pi i(m+\alpha)x}.
\]
Both families are orthonormal bases.  Let $P_0$ be the projection onto
the closed span of $\{e_n:n\le0\}$ and $P_\alpha$ the projection onto
the closed span of $\{f_m^{(\alpha)}:m\le0\}$.  The shift
$\alpha=k/4$ is exactly the momentum displacement
\[
 \Delta p=\frac{2\pi k}{4N_x}
\]
relative to the spacing $2\pi/N_x$.

\begin{lemma}[Exact overlap]\label{lem:overlap}
For all $m,n\in\Z$,
\[
 \left|\ip{e_n}{f_m^{(\alpha)}}\right|^2
 =\frac{\sin^2(\pi\alpha)}
 {\pi^2(m-n+\alpha)^2}.
\tag{5.1}\label{eq:overlap}
\]
\end{lemma}

\begin{proof}
Direct integration gives
\[
\begin{aligned}
\ip{e_n}{f_m^{(\alpha)}}
&=\int_0^1e^{2\pi i(m-n+\alpha)x}\,dx\\
&=\frac{e^{2\pi i(m-n+\alpha)}-1}
 {2\pi i(m-n+\alpha)}\\
&=e^{\pi i(m-n+\alpha)}
\frac{\sin(\pi(m-n+\alpha))}
 {\pi(m-n+\alpha)}.
\end{aligned}
\]
Since $m-n$ is an integer,
$|\sin(\pi(m-n+\alpha))|=|\sin(\pi\alpha)|$.
Taking the squared modulus proves \eqref{eq:overlap}.
\end{proof}

\begin{theorem}[One-edge Hilbert--Schmidt asymptotic]
\label{thm:oneedge}
Let $S_N(\alpha)$ be the contribution to
$\norm{P_0-P_\alpha}_2^2$ obtained by restricting both occupied and
unoccupied indices to absolute value at most $N$.  Then
\[
 S_N(\alpha)
 =\frac{2\sin^2(\pi\alpha)}{\pi^2}\log N+O_\alpha(1).
\tag{5.2}\label{eq:one-edge-log}
\]
In particular, $P_0-P_\alpha$ is not Hilbert--Schmidt.
\end{theorem}

\begin{proof}
For any two orthogonal projections $P,Q$ for which the terms are
defined,
\[
 (P-Q)^2=P(\id-Q)P+(\id-P)Q(\id-P).
\]
Taking the trace in orthonormal bases gives
\[
 \norm{P-Q}_2^2
 =\Tr(P(\id-Q))+\Tr(Q(\id-P)).
\tag{5.3}\label{eq:projection-distance}
\]
For the first term in \eqref{eq:projection-distance}, occupied
$e_n$ have $n\le0$ and unoccupied $f_m^{(\alpha)}$ have $m\ge1$.
For the second, occupied $f_m^{(\alpha)}$ have $m\le0$ and
unoccupied $e_n$ have $n\ge1$.  By Lemma~\ref{lem:overlap}, the
symmetrically truncated sum is
\[
 S_N(\alpha)
 =\frac{\sin^2(\pi\alpha)}{\pi^2}
 \sum_{d=1}^{2N}\mu_N(d)
 \left(\frac1{(d+\alpha)^2}+\frac1{(d-\alpha)^2}\right),
\tag{5.4}\label{eq:finite-sum}
\]
where
\[
 \mu_N(d)=
 \#\{(m,n):1\le m\le N,\ -N\le n\le0,\ m-n=d\}.
\]
Counting the integer points on each diagonal gives
\[
 \mu_N(d)=
 \begin{cases}
 d,&1\le d\le N,\\
 2N-d+1,&N<d\le2N.
 \end{cases}
\tag{5.5}\label{eq:multiplicity}
\]
The possible one-unit convention change at $d=N+1$ affects only a
bounded term and not the asymptotic; \eqref{eq:multiplicity} corresponds
to the displayed index ranges.

For $d\ge1$,
\[
\begin{aligned}
d\left(\frac1{(d+\alpha)^2}+\frac1{(d-\alpha)^2}\right)
&=\frac1d\left(
\frac1{(1+\alpha/d)^2}
+\frac1{(1-\alpha/d)^2}\right)\\
&=\frac2d+O_\alpha(d^{-3}).
\end{aligned}
\tag{5.6}\label{eq:summand-expand}
\]
The absence of a $d^{-2}$ term follows because the terms odd in
$\alpha/d$ cancel.  Therefore the part $1\le d\le N$ of
\eqref{eq:finite-sum} equals
\[
 2\sum_{d=1}^N\frac1d+O_\alpha\left(\sum_{d=1}^\infty d^{-3}\right)
 =2\log N+O_\alpha(1).
\tag{5.7}\label{eq:harmonic}
\]
For the tail $N<d\le2N$, use $d-\alpha\ge d/2$ to obtain
\[
\begin{aligned}
\sum_{d=N+1}^{2N}\mu_N(d)
\left(\frac1{(d+\alpha)^2}+\frac1{(d-\alpha)^2}\right)
&\le
\sum_{d=N+1}^{2N}(2N-d+1)\frac8{d^2}\\
&\le \frac8{N^2}\sum_{r=1}^{N}r
\le 4+\frac4N.
\end{aligned}
\tag{5.8}\label{eq:tail}
\]
Thus the tail is uniformly bounded.  Substitution of
\eqref{eq:harmonic} and \eqref{eq:tail} into \eqref{eq:finite-sum}
proves \eqref{eq:one-edge-log}.  Since the nonnegative truncations
diverge, the full sum in \eqref{eq:projection-distance} is infinite,
so $P_0-P_\alpha$ is not Hilbert--Schmidt.
\end{proof}

\subsection{Two edges and comparison with the QWZ measurement}

\begin{corollary}[Two-edge quarter twist]\label{cor:twoedge}
For two independent chiral edge branches, one at each boundary of a
cylinder, the direct-sum polarizations satisfy
\[
 \norm{P_0^{\mathrm{edge}}-P_\alpha^{\mathrm{edge}}}_2^2
 =\frac{4\sin^2(\pi\alpha)}{\pi^2}\log N+O_\alpha(1).
\tag{5.9}\label{eq:two-edge-log}
\]
For the adjacent $\mu_4$ sectors, $\alpha=1/4$, hence
\[
 \norm{P_0^{\mathrm{edge}}-P_{1/4}^{\mathrm{edge}}}_2^2
 =\frac2{\pi^2}\log N+O(1),
\qquad
\frac2{\pi^2}=0.2026423672846755\ldots .
\tag{5.10}\label{eq:quarter-coefficient}
\]
The difference from the measured slope $0.203$ is
$0.0003576327\ldots$, or approximately $0.1762\%$ of $0.203$.
\end{corollary}

\begin{proof}
Hilbert--Schmidt norms square add under orthogonal direct sums.  The two
edges have opposite propagation orientation, but the overlap formula
depends on the squared modulus and is unchanged.  Doubling
\eqref{eq:one-edge-log} proves \eqref{eq:two-edge-log}.  Since
$\sin^2(\pi/4)=1/2$, the coefficient becomes $2/\pi^2$.  The numerical
comparison is direct arithmetic.
\end{proof}

\begin{lemma}[Stability under a uniformly Hilbert--Schmidt remainder]
\label{lem:stability}
Suppose finite-volume differences $D_N$ admit
\[
 D_N=D_N^{\mathrm{edge}}+R_N,\qquad
 \sup_N\norm{R_N}_2<\infty,
\tag{5.11}\label{eq:remainder}
\]
and
\[
 \norm{D_N^{\mathrm{edge}}}_2^2=c\log N+O(1)
\quad(c>0).
\]
Then
\[
 \norm{D_N}_2^2=c\log N+O(\sqrt{\log N}),
\tag{5.12}\label{eq:stable-log}
\]
so $\norm{D_N}_2\to\infty$ and the leading logarithmic coefficient is
still $c$.
\end{lemma}

\begin{proof}
Expand the square:
\[
 \norm{D_N}_2^2
 =\norm{D_N^{\mathrm{edge}}}_2^2
 +2\operatorname{Re}\Tr((D_N^{\mathrm{edge}})^*R_N)
 +\norm{R_N}_2^2.
\]
Hilbert--Schmidt Cauchy--Schwarz gives
\[
 \left|\Tr((D_N^{\mathrm{edge}})^*R_N)\right|
 \le\norm{D_N^{\mathrm{edge}}}_2\norm{R_N}_2
 =O(\sqrt{\log N}).
\]
The last term is $O(1)$ by \eqref{eq:remainder}.  This proves
\eqref{eq:stable-log}.
\end{proof}

\begin{theorem}[Global non-implementability in the proved edge class]
\label{thm:nonimpl}
Let $P_0$ and $P_{1/4}$ be the infinite-volume pure quasi-free
polarizations of either the exact two-edge model of
Corollary~\ref{cor:twoedge}, or a lattice collar for which the
uniform-remainder decomposition \eqref{eq:remainder} holds along an
exhaustion.  Then the two quasi-free representations are not unitarily
equivalent.  Equivalently, no global Fock-space Bogoliubov implementer
can carry the neutral polarization to the quarter-twisted polarization.
\end{theorem}

\begin{proof}
For the exact edge model, Theorem~\ref{thm:oneedge} and
Corollary~\ref{cor:twoedge} show that the nonnegative finite-rank
truncations of $\norm{P_0-P_{1/4}}_2^2$ tend to infinity.  Therefore
$P_0-P_{1/4}\notin\mathcal L^2$.  Under
\eqref{eq:remainder}, Lemma~\ref{lem:stability} gives the same conclusion.

Assume for contradiction that a global Bogoliubov implementer exists.
Let $R$ be the underlying one-particle orthogonal transformation, so
$P_{1/4}=RP_0R^*$.  The Shale--Stinespring criterion,
Theorem~\ref{thm:ss}, would imply
\[
 [P_0,R]\in\mathcal L^2.
\]
But
\[
 [P_0,R]R^*=P_0-RP_0R^*=P_0-P_{1/4},
\]
and multiplication by $R^*$ preserves the Hilbert--Schmidt class.
Hence $P_0-P_{1/4}$ would be Hilbert--Schmidt, contradicting the
divergence just proved.  The pure-state Araki criterion quoted after
Theorem~\ref{thm:ps} gives the equivalent representation-theoretic
formulation.
\end{proof}

\begin{remark}[Gapped control]
If instead the global projection difference converges in
Hilbert--Schmidt norm, Shale--Stinespring permits a global implementer.
The measured trivial-phase plateau near $0.546$ is consistent with this
alternative, but a finite plateau is not a proof of convergence.
\end{remark}

\section{Combined strongest theorem}

\begin{theorem}[Strongest theorem proved from the presently isolated
mathematical hypotheses]\label{thm:main}
Consider a sequence of finite quadratic fermion systems with four
$\mu_4$-twisted sectors and a nested family of fixed finite windows.
Assume the following precise hypotheses.
\begin{enumerate}[label=(H\arabic*)]
\item On each fixed window, the restricted covariances
$C_{N,w}^{(k)}$ converge in trace norm.  For the adjacent sectors
$k=0,1$, their limiting Araki--Wyss projections have distance strictly
less than one.  In particular, this holds when the two limiting window
covariances coincide.
\item The finite observable algebras form an equivariant unital tower,
and their clock unitaries obey
$u_N^4=\id$ and $\iota_N(u_N)=u_{N+1}$.
\item The grade-one class is labelled by the frozen lattice element
$[\lambda]$ with $[\lambda]^4=[0]$ and $[\lambda]^2=[v]$.
\item For the global topological conclusion, the infinite-volume
polarization difference is the two-edge quarter-twist model of
Section~\ref{sec:global}, up to a remainder uniformly bounded in
Hilbert--Schmidt norm along the exhaustion.
\end{enumerate}
Then:
\begin{enumerate}[label=\textup{(\alph*)}]
\item On every fixed window, the adjacent sector change has a canonical
Bogoliubov implementer in the purified standard representation, unique
up to phase, with the explicit bound
\[
\norm{[\widehat P_{C_{N,w}^{(0)}},R_{N,w}]}_2
\le2\sqrt{2\norm{C_{N,w}^{(0)}-C_{N,w}^{(1)}}_1}.
\]
\item Along every subsequence on which the phase is chosen coherently,
these implementers converge in operator norm, hence strongly on the
finite-particle core; estimate \eqref{eq:nparticle} holds for every
$n$-particle vector.
\item The inductive-limit extension is
\[
\A\rtimes\mathbb Z_4\cong\A\otimes C^*(\mathbb Z_4),
\]
the homogeneous subspaces have $\mathbb Z_4$ fusion,
$\Psi_\lambda^4\in\A$, and $\Psi_\lambda^2$ has the parity grade
$[v]$.  The coefficient expectation is the unique dual-invariant
expectation onto $\A$, and the group average is the unique
trace-preserving expectation onto $\A^{\mathbb Z_4}$.
\item In the topological two-edge class no global Bogoliubov implementer
exists.  The obstruction has the exact leading law
\[
\norm{P_0-P_{1/4}}_2^2
=\frac2{\pi^2}\log N+O(\sqrt{\log N}).
\]
\end{enumerate}
\end{theorem}

\begin{proof}
Parts (a) and (b) are Theorem~\ref{thm:local}, with all hypotheses
verified in its proof.  Part (c) is Theorem~\ref{thm:crossed}; hypothesis
(H3) supplies only the identification of the abstract grades with the
frozen lattice labels.  Part (d) follows from
Corollary~\ref{cor:twoedge}, Lemma~\ref{lem:stability}, and
Theorem~\ref{thm:nonimpl}.  No further implication is used.
\end{proof}

\section{What this theorem does and does not close}

\subsection{Relation to G1--G4}

\begin{description}[style=nextline]
\item[G1: convergence on the finite-particle core with energy bounds.]
Theorem~\ref{thm:local} proves strong convergence on each
\emph{fixed purified finite window}.  It does not prove convergence of
a nontrivial charged intertwiner on one common infinite-volume edge Fock
space.  The measured inequalities with $C=2.86$ concern only
$n=1,\ldots,4$ and finite $N_x$; they do not imply
\[
\norm{H_{\mathrm{edge}}^n\Psi_\lambda\Omega}
\le C_0^{n+1}(n!)^\beta\qquad\text{for every }n\ge0.
\]
Thus G1 remains open in its contract meaning.

\item[G2: locality relative to the bilinear net.]
No commutator estimate at spacelike separation, no Haag--Kastler
localisation, and no relative-locality theorem follows from covariance
convergence.  G2 is untouched.

\item[G3: braiding and monodromy.]
The exact arithmetic $h_\lambda=1$ is the correct trivial
self-monodromy precondition.  The measured Berry phase $0$ with spectral
flow $1$ is consistent with it.  Neither fact computes the DHR braiding
operator or proves the relative exchange relation of a charged field.
G3 is therefore not closed.

\item[G4: Q-system closure.]
Section~\ref{sec:tower} proves the $\mathbb Z_4$-graded crossed-product
algebra, its naturality in the finite tower, and the precisely normalised
expectations.  This is a complete $C^*$-algebraic statement for the inner
clock tower.  It does not identify that tower with the local
simple-current Q-system of the $\mathfrak{so}(16)_1$ conformal net.
Hence it is the exact finite/inductive algebraic shadow of G4, not the
net-level G4 demanded by the contract.
\end{description}

\subsection{The MMST/Osborne--Stottmeister leg}

T.~J.~Osborne and A.~Stottmeister,
\emph{Conformal field theory from lattice fermions},
Communications in Mathematical Physics \textbf{398} (2023), 219--283,
\href{https://arxiv.org/abs/2107.13834}{arXiv:2107.13834},
prove strong convergence of the lattice fermion-bilinear currents on a
finite-particle core (Theorem~5.1 for the displayed $u(1)$ current, with
the nonabelian bilinear extension by the Cauchy--Schwarz reduction around
their Eq.~(229)), convergence of the Koo--Saleur/Virasoro generators
(Theorems~4.11 and 4.16), and correlation-function control
(Theorems~6.1 and 6.3).  Their error estimates around Eqs.~(205) and
(207) are polynomial in lattice spacing.

Those theorems concern observables that are fermion bilinears.  The
$128$ weight-one spinor currents needed for
$\mathfrak{so}(16)_1\to(E_8)_1$ are not fermion bilinears, and the twist
field is outside that proved class.  Therefore the statement additionally
needed here is not another application of Shale--Stinespring.  It is an
identification theorem saying that one specifically normalised limit of
the lattice twist intertwiners is a relatively local DHR charged field
whose $\mathbb Z_4$ Q-system extension is the holomorphic $(E_8)_1$ net.
That assertion belongs to conformal-net/simple-current extension theory,
not to the quasi-free category: a spin field changes sectors and is not a
polynomial fermion bilinear.

\subsection{Status of the measured QWZ logarithm}

The coefficient $0.203$ has a sharp analytic explanation:
each chiral boundary contributes
\[
\frac{2\sin^2(\pi/4)}{\pi^2}=\frac1{\pi^2},
\]
and the two boundaries contribute $2/\pi^2=0.202642367\ldots$.
The measured intercept $3.605$ is nonuniversal and receives bulk,
ultraviolet, and finite-width contributions.  To turn the coefficient
match into a proof for the QWZ collar, one must establish the decomposition
\eqref{eq:remainder} uniformly in $N_x$.  That estimate is not contained
in the frozen probes.

\section{Remaining gap}

\begin{center}
\setlength{\fboxsep}{10pt}
\fcolorbox{red!65!black}{red!3}{
\begin{minipage}{0.91\textwidth}
\textbf{Remaining gap (exact statement).}
For the $M=1$, $N_y=8$ QWZ collar with quarter-twisted sectors, construct
isometric embeddings of all finite-volume one-particle spaces into one
edge Hilbert space and prove simultaneously:
\[
C_N^{(0)}-C_N^{(1)}
=D_{N,\mathrm{top}}^{(1/4)}
\oplus D_{N,\mathrm{bottom}}^{(1/4)}+R_N,
\qquad
\sup_N\norm{R_N}_{\mathrm{HS}}<\infty;
\]
and there exist phase-coherent lattice intertwiners
$\Psi_{\lambda,N}$ converging on one common finite-energy core to a
nonzero operator $\Psi_\lambda$ such that, for constants
$C_0<\infty$ and $\beta<\infty$ independent of $n$,
\[
\norm{H_{\mathrm{edge}}^n\Psi_\lambda\Omega}
\le C_0^{\,n+1}(n!)^\beta\quad(n\ge0),
\]
$\Psi_\lambda$ is relatively local with respect to the
$\mathfrak{so}(16)_1$ bilinear net, its DHR exchange operators realise
the order-four class with $h_\lambda=1$ and
$[\lambda]^2=[v]$, and the local extension generated by
$\{\mathcal A_{\mathfrak{so}(16)_1},\Psi_\lambda\}$ is isomorphic, as a
conformal net with vacuum and stress tensor, to the holomorphic
$(E_8)_1$ net.

\medskip
The first line upgrades the finite logarithmic fit to the proved Shale
obstruction for the actual collar.  The remaining lines are the
MMST/Osborne--Stottmeister extension beyond fermion bilinears and supply
G1--G4 at net level.  Neither assertion follows from the measured
windowed Hilbert--Schmidt Cauchy rate, three phase-fixed matrix elements,
or four finite energy moments.
\end{minipage}}
\end{center}

\section*{Conclusion}
\addcontentsline{toc}{section}{Conclusion}

The quasi-free conclusions that can be proved today are complete:
finite-window standard-form implementers with an explicit Shale bound,
their coherent strong convergence, the natural inner-$\mathbb Z_4$
crossed-product tower with correctly qualified unique expectations, and
the universal two-edge quarter-twist obstruction with coefficient
$2/\pi^2$.  The data strongly match this structure, but they do not yet
construct the net-level field $\Psi_\lambda$.  The remaining box states
exactly the theorem that would be required to do so.

\begin{thebibliography}{9}

\bibitem{Araki1970}
H.~Araki,
\emph{On quasifree states of CAR and Bogoliubov automorphisms},
Publ. Res. Inst. Math. Sci. \textbf{6} (1970/71), 385--442.

\bibitem{OsborneStottmeister2023}
T.~J.~Osborne and A.~Stottmeister,
\emph{Conformal field theory from lattice fermions},
Commun. Math. Phys. \textbf{398} (2023), 219--283,
\href{https://arxiv.org/abs/2107.13834}{arXiv:2107.13834}.

\bibitem{PowersStormer1970}
R.~T.~Powers and E.~St{\o}rmer,
\emph{Free states of the canonical anticommutation relations},
Commun. Math. Phys. \textbf{16} (1970), 1--33.

\bibitem{Ruijsenaars1978}
S.~N.~M.~Ruijsenaars,
\emph{On Bogoliubov transformations. II. The general case},
Ann. Physics \textbf{116} (1978), 105--134.

\bibitem{ShaleStinespring1964}
D.~Shale and W.~F.~Stinespring,
\emph{States of the Clifford algebra},
Ann. of Math. (2) \textbf{80} (1964), 365--381.

\end{thebibliography}

\end{document}
